\chapter{Copper (Blood)}

\section*{What Is This Test?}
The serum copper test measures the concentration of copper, an essential trace element, in the blood. Copper is a cofactor for numerous enzymes, including cytochrome c oxidase, superoxide dismutase, dopamine beta-hydroxylase, lysyl oxidase, and ceruloplasmin (the main copper transport protein). Approximately 90--95\% of serum copper is bound to ceruloplasmin; the remainder is loosely bound to albumin and other proteins (the ``free'' or non-ceruloplasmin copper fraction). Copper is absorbed in the proximal small intestine, transported to the liver, and incorporated into ceruloplasmin, which is then secreted into the circulation; excess copper is excreted in bile. Serum copper is measured to evaluate copper deficiency (which can cause anemia, neutropenia, and neurologic disease), copper excess, and, most importantly, as part of the diagnostic workup of Wilson disease, a genetic disorder of hepatic copper excretion that causes copper accumulation in the liver, brain, and cornea. Serum copper is interpreted together with serum ceruloplasmin, 24-hour urine copper, and (when indicated) free copper and hepatic copper measurement.

\section*{Alternative Names}
Serum Copper; Copper Level; Cu; Copper, Blood; Plasma Copper

\section*{Principle}
Serum copper is measured by atomic absorption spectroscopy (AAS) or, increasingly, by inductively coupled plasma mass spectrometry (ICP-MS), both of which provide parts-per-billion sensitivity. The specimen is collected in a trace-element-free (royal blue-top) tube to avoid contamination, because copper is ubiquitous in the environment and in ordinary collection tubes. Total serum copper is the sum of ceruloplasmin-bound copper (the large majority) and non-ceruloplasmin (free) copper. Ceruloplasmin binds six copper atoms per molecule and is an acute-phase reactant; therefore serum copper rises and falls with ceruloplasmin, which complicates interpretation in inflammation, pregnancy, and estrogen therapy. The ``free'' (non-ceruloplasmin) copper is calculated as: free copper = total serum copper $-$ (ceruloplasmin [mg/dL] $\times$ 3.15 [mcg Cu per mg ceruloplasmin]). In Wilson disease, total serum copper is low (because ceruloplasmin synthesis is reduced) and free copper is elevated (because the impaired biliary excretion causes copper to accumulate and overflow into the non-ceruloplasmin compartment). Copper is also measurable in 24-hour urine (markedly elevated in Wilson disease due to the overflow of free copper) and directly in liver tissue (the diagnostic gold standard).

\section*{History}
Copper was shown to be an essential nutrient in 1928 by Elvehjem and Hart, who demonstrated that rats fed an iron-supplemented milk diet still developed anemia unless copper was added, establishing copper as essential for hemoglobin formation. Wilson disease was first described in 1912 by the British neurologist Samuel Alexander Kinnier Wilson as ``progressive lenticular degeneration.'' The role of copper in Wilson disease was established over the following decades: in 1948, the presence of Kayser-Fleischer rings was linked to copper deposition, and in 1952, Scheinberg and Gitlin demonstrated low serum ceruloplasmin in Wilson disease, providing the basis for the serum copper and ceruloplasmin tests. In 1956, Walshe introduced penicillamine, the first effective copper-chelating therapy, and later zinc acetate for maintenance. The Menkes disease (copper deficiency) syndrome was described in 1962 by John Menkes. The Wilson disease gene, ATP7B, was cloned in 1993 by three groups, enabling genetic diagnosis and the modern era of population and family screening.

\section*{Why This Test Is Ordered}
\begin{itemize}
  \item \textbf{Evaluation of suspected Wilson disease:} In any patient (typically children and young adults) with unexplained liver disease (hepatitis, cirrhosis, acute liver failure), neurologic symptoms (tremor, dysarthria, dystonia, parkinsonism), psychiatric symptoms, or Kayser-Fleischer rings
  \item \textbf{Screening of first-degree relatives of a patient with Wilson disease:} Because the condition is autosomal recessive
  \item \textbf{Evaluation of copper deficiency:} In patients with unexplained anemia, neutropenia, or myeloneuropathy, particularly after bariatric surgery or upper gastrointestinal surgery, or with excessive zinc supplementation or malabsorption
  \item \textbf{Evaluation of Menkes disease (infants):} Characterized by kinky hair, failure to thrive, and progressive neurodegeneration with severe copper deficiency
  \item \textbf{Evaluation of copper toxicity:} Accidental ingestion, contaminated water, chronic cholestatic liver disease with copper accumulation (primary biliary cholangitis)
  \item \textbf{Monitoring of Wilson disease treatment:} With chelators (penicillamine, trientine) and zinc therapy
  \item \textbf{Evaluation of unexplained liver disease in children and young adults:} As part of the metabolic workup
\end{itemize}

\section*{Is This a Primary or Secondary Test}
Serum copper is a secondary (diagnostic) test. For Wilson disease, no single test is diagnostic; the diagnosis rests on a combination of serum ceruloplasmin, serum copper (total and free), 24-hour urinary copper, and often hepatic copper or genetic testing (ATP7B). Serum copper is always interpreted with ceruloplasmin.

\section*{How This Test Is Performed}
A venous blood sample is collected in a royal blue-top (trace-element-free) tube, typically 3--5 mL, using strict technique to avoid environmental contamination (skin preparation, avoidance of hemolysis, and metal-free tubes). The sample may be drawn at any time of day; fasting is not required. The serum is analyzed by atomic absorption spectroscopy or ICP-MS, with results available within 1--3 days. For a Wilson disease workup, a 24-hour urine collection for copper and serum ceruloplasmin are drawn concurrently. The urine collection must be in an acid-washed, metal-free container, and the laboratory result is expressed as mcg per 24 hours.

\section*{What Preparation Is Needed}
No fasting is required, although some laboratories recommend a morning draw. The patient should inform the healthcare provider of: all supplements (copper, zinc, multivitamins --- zinc interferes with copper absorption), estrogen or oral contraceptive use (which raises ceruloplasmin and serum copper), pregnancy, recent surgery, and any known liver disease. Medications such as penicillamine and trientine must be disclosed, as they lower serum copper. Because zinc supplementation is the standard maintenance therapy for Wilson disease and lowers copper, patients on treatment should not stop it without medical advice.

\section*{Contraindications and Precautions}
There are no absolute contraindications to the blood draw. Important interpretive cautions: (1) serum copper is an unreliable single test for Wilson disease --- it may be normal in up to one-third of patients, particularly early in the disease or during acute liver failure (when massive hepatocellular release can raise it); (2) low ceruloplasmin (and thus low total copper) also occurs in non-Wilson conditions such as severe protein-losing states, end-stage liver disease of any cause, nephrotic syndrome, and Menkes disease; (3) ceruloplasmin is an acute-phase reactant, so inflammation, pregnancy, and estrogen therapy raise both ceruloplasmin and serum copper, potentially masking a low level; (4) hemolyzed specimens falsely elevate serum copper (red cells are copper-rich); (5) contamination from skin, needles, or standard tubes causes false elevations --- trace-element collection is mandatory; (6) free (non-ceruloplasmin) copper greater than approximately 15 mcg/dL is a useful but calculated value and is not measured directly; (7) a low serum copper in a patient not on therapy raises concern for copper deficiency, which must be corrected before starting or continuing zinc.

\section*{Risks}
Risks are those of routine venipuncture: mild pain, bruising, small hematoma at the collection site, lightheadedness or vasovagal syncope, and extremely rarely, infection or nerve injury. The 24-hour urine collection is noninvasive and carries no risk.

\section*{What Happens After the Test}
If Wilson disease is suspected and serum copper and ceruloplasmin are low with elevated 24-hour urinary copper and elevated free copper, the diagnosis is confirmed (with hepatic copper or ATP7B genetic testing when needed) and treatment with a chelator (penicillamine or trientine) or zinc is initiated, with lifelong monitoring. A slit-lamp examination is performed to look for Kayser-Fleischer rings. If copper deficiency is found (low serum copper with low ceruloplasmin, typically after bariatric surgery or with excessive zinc), oral or intravenous copper supplementation is given, and the underlying cause (zinc overuse, malabsorption) is addressed. If copper excess without Wilson disease is present (cholestatic liver disease, contamination), the cause is investigated. The healthcare provider interprets serum copper only in the context of ceruloplasmin, urinary copper, and the clinical picture.

\section*{Understanding Results}

\subsection*{Below Normal (Copper Deficiency)}
\begin{itemize}
  \item \textbf{Wilson disease:} Low total serum copper and low ceruloplasmin with high free copper and high urinary copper; requires treatment
  \item \textbf{Copper deficiency (acquired):} Low serum copper after bariatric or upper gastrointestinal surgery, prolonged parenteral nutrition without copper, excessive zinc supplementation (zinc induces intestinal metallothionein, blocking copper absorption), malabsorption, or malnutrition; causes anemia (often with neutropenia), and myeloneuropathy with sensory ataxia
  \item \textbf{Menkes disease (infants):} Severe copper deficiency with kinky hair, hypotonia, seizures, and connective tissue abnormalities
  \item \textbf{Nephrotic syndrome or severe protein loss:} Low ceruloplasmin with consequent low serum copper
\end{itemize}

\subsection*{Above Normal (Copper Excess)}
\begin{itemize}
  \item \textbf{Acute copper toxicity:} Ingestion or accidental exposure; causes gastroenteritis, hemolysis, hepatic and renal failure
  \item \textbf{Cholestatic liver disease:} Primary biliary cholangitis and other cholestatic syndromes accumulate copper in the liver with variable serum levels
  \item \textbf{Physiological and drug-related elevation:} Pregnancy, estrogen therapy and oral contraceptives (raise ceruloplasmin), and inflammation (acute-phase response)
  \item \textbf{Wilson disease with fulminant hepatic failure:} Massive hepatocellular copper release can transiently raise serum copper
  \item \textbf{Contamination or hemolysis:} False elevation from improper collection or red cell lysis
\end{itemize}

\section*{Normal Results}
\begin{itemize}
  \item \textbf{Serum copper (adults):} 70--140 mcg/dL (11--22 $\mu$mol/L)
  \item \textbf{Serum copper (children):} Approximately 80--150 mcg/dL (reference ranges vary with age)
  \item \textbf{Serum copper (pregnancy and oral contraceptive use):} Elevated, up to 150--250 mcg/dL due to increased ceruloplasmin
  \item \textbf{Free (non-ceruloplasmin) copper:} Normally $<$10--15 mcg/dL; values above approximately 15--25 mcg/dL support Wilson disease
  \item \textbf{24-hour urinary copper:} Normally $<$40--60 mcg/24 hours (some laboratories use $<$100 mcg/24 hours); Wilson disease typically shows $>$100 mcg/24 hours at diagnosis (and often $>$1,000 mcg/24 hours)
  \item \textbf{Serum ceruloplasmin:} 20--40 mg/dL (low in Wilson disease, usually $<$20 mg/dL)
\end{itemize}

\section*{Follow-up Tests}
\begin{itemize}
  \item \textbf{Serum ceruloplasmin:} Always interpreted with serum copper; low in Wilson disease
  \item \textbf{24-hour urine copper:} Elevated in untreated Wilson disease; also used to monitor chelation therapy (target 200--500 mcg/24 hours)
  \item \textbf{Free (non-ceruloplasmin) copper:} Calculated from total copper and ceruloplasmin to support the diagnosis
  \item \textbf{Liver function tests and coagulation:} To assess hepatic involvement in Wilson disease
  \item \textbf{Hepatic copper quantification (liver biopsy):} The gold standard, with $>$250 mcg/g dry weight supporting Wilson disease
  \item \textbf{ATP7B genetic testing:} Confirms the diagnosis and enables family screening
  \item \textbf{Slit-lamp examination:} For Kayser-Fleischer rings (corneal copper deposition)
  \item \textbf{Complete blood count and neurologic examination:} To evaluate hematologic and neurologic manifestations of deficiency
  \item \textbf{Urine zinc:} When copper deficiency from zinc excess is suspected
\end{itemize}

\section*{References}
\begin{enumerate}
  \item Schilsky ML, Roberts EA, Bronstein JM, et al. A multidisciplinary approach to the diagnosis and management of Wilson disease: 2022 practice guidance on Wilson disease from the American Association for the Study of Liver Diseases. Hepatology. 2023;77(4):1428--1455.
  \item Roberts EA, Schilsky ML. Diagnosis and treatment of Wilson disease: an update. Hepatology. 2008;47(6):2089--2111.
  \item Burtis CA, Ashwood ER, Bruns DE. Tietz Textbook of Clinical Chemistry and Molecular Diagnostics. 7th ed. Elsevier; 2023.
  \item Fischbach FT, Dunning MB. A Manual of Laboratory and Diagnostic Tests. 11th ed. Wolters Kluwer; 2023.
  \item Scheinberg IH, Gitlin D. Deficiency of ceruloplasmin in patients with hepatolenticular degeneration (Wilson's disease). Science. 1952;116(3018):484--485.
  \item Kumar N. Copper deficiency myelopathy (human swayback). Mayo Clin Proc. 2006;81(10):1371--1384.
  \item Menkes JH, Alter M, Steigleder GK, Weakley DR, Sung JH. A sex-linked recessive disorder with retardation of growth, peculiar hair, and focal cerebral and cerebellar degeneration. Pediatrics. 1962;29:764--779.
  \item Walshe JM. Penicillamine, a new oral therapy for Wilson's disease. Am J Med. 1956;21(4):487--495.
\end{enumerate}

\clearpage
